% TOP_PROBLEM: {"id": 30006742, "title": "Polynomial pivot rule for the simplex method", "category_id": 15, "category": {"id": 15, "name": "computer_science", "display_name": "Computer Science", "description": "Computational complexity, algorithms, and theoretical CS.", "slug": "computer-science", "order_index": 15, "created_at": "2026-07-31T15:26:25.671Z"}, "status": "open"}
% FIRST_POSTED: 2026-09-27
% !TeX program = lualatex
% Compile twice: lualatex 30006742.tex
% All references and prose are embedded; no BibTeX database or external inputs.
\documentclass[11pt,a4paper]{article}
\usepackage[margin=24mm,headheight=14pt]{geometry}
\usepackage{fontspec}
\setmainfont{Latin Modern Roman}
\setsansfont{Latin Modern Sans}
\setmonofont{Latin Modern Mono}
\usepackage{amsmath,amssymb,mathtools}
\usepackage{unicode-math}
\setmathfont{Latin Modern Math}
\usepackage{microtype}
\usepackage{enumitem,xurl,needspace}
\usepackage[dvipsnames]{xcolor}
\usepackage{hyperref}
\hypersetup{unicode=true,colorlinks=true,linkcolor=MidnightBlue,urlcolor=MidnightBlue,
 pdftitle={500 Mathematical Problem Statements},pdfauthor={Source-annotated compilation},bookmarksnumbered=true}
\usepackage{bookmark}
\usepackage{tocloft}
\setlength{\cftsecnumwidth}{3em}
\cftsetpnumwidth{2.5em}
\cftsetrmarg{4em}
\usepackage{titlesec}
\titleformat{\section}{\Large\bfseries\raggedright}{\thesection}{0.7em}{}
\usepackage{fancyhdr}
\pagestyle{fancy}\fancyhf{}
\fancyhead[L]{\small\sffamily Mathematical problem statements}
\fancyhead[R]{\small Source edition 22}
\fancyfoot[C]{\thepage}
\setlength{\parindent}{0pt}
\setlength{\parskip}{5pt plus 1pt}
\setlength{\emergencystretch}{3em}
\setcounter{tocdepth}{1}
\setcounter{secnumdepth}{1}
\setlist[itemize]{leftmargin=3em,itemsep=5pt,topsep=4pt,parsep=0pt}
\allowdisplaybreaks
\newcommand{\N}{\mathbb{N}}
\newcommand{\Z}{\mathbb{Z}}
\newcommand{\Q}{\mathbb{Q}}
\newcommand{\R}{\mathbb{R}}
\newcommand{\C}{\mathbb{C}}
\newcommand{\F}{\mathbb{F}}
\newcommand{\A}{\mathbb{A}}
\newcommand{\PP}{\mathbb P}
\newcommand{\E}{\mathbb{E}}
\newcommand{\eps}{\varepsilon}
\newcommand{\dd}{\,\mathrm d}
\newcommand{\sref}[2]{\hyperlink{source:#1:#2}{[S#2]}}
\newcommand{\eref}[2]{\hyperlink{extra:#1:#2}{[E#2]}}
% Turn the next switch off for a shorter reading copy. The quotations preserve
% every exactTarget field, including qualifications omitted from a short summary.
\newif\ifshowcatalogue
\showcataloguetrue
\newcommand{\cataloguescope}[1]{\ifshowcatalogue
 \paragraph{Exact scope in the supplied catalogue.}
 \begin{quote}\small #1\end{quote}\fi}
\hypersetup{pdftitle={Problem 30006742 - Polynomial simplex pivot count},pdfauthor={Open Problems Math research notebook}}
\fancyhead[L]{\small\sffamily Open Problems Math \textbar\ Research notebook}
\fancyhead[R]{\small\sffamily 131 \textbar\ 27 September 2026}
\numberwithin{equation}{section}
\begin{document}
\setcounter{section}{0}
% ================================================================
% problemId: 30006742
% Canonical problem ID: 30006742
% exactTarget SHA-256: 84962be046fea3ae2dd52dfba4ffb327b67213cc9288d69c41ba9839897d1fda
\clearpage
\section*{\texorpdfstring{Polynomial pivot rule for the simplex method}{Polynomial pivot rule for the simplex method}}\label{30006742}
\subsection{Definitions and mathematical statement}
In standard form a rational linear program has $Ax=b$, $x\ge0$, and a linear objective. A feasible basis selects an invertible basic submatrix and yields a feasible basic solution. A simplex pivot exchanges one basic and one nonbasic variable according to the rule while maintaining the algorithm's feasibility and objective conventions. For bounded nondegenerate inputs, seek a rule with
\[
 \sup_{\text{allowed LPs and feasible bases}}\mathbb E[\text{pivots to an optimum}]
 \le p(m,n)
\]
for one polynomial independent of coefficient bit lengths. For a deterministic rule omit expectation. The catalogue requests a pivot-count bound and does not separately impose an efficient per-pivot selection model.
\par\smallskip\textit{Formulation and concept references:} \sref{30006742}{1}.
\cataloguescope{Construct a deterministic or randomized simplex pivot rule and a polynomial p such that, for every bounded nondegenerate rational linear program with m equality constraints and n variables and every feasible starting basis, the rule reaches an optimal basis in at most p(m,n) pivots, in expectation when the rule is randomized; or prove that no such rule exists.}
\subsection{Short English statement}
Is there a simplex pivot rule that always reaches an optimum after polynomially many pivots, independent of coefficient magnitudes?
% BEGIN RESEARCH ADDITION 131
\subsection{Research attempt}

\paragraph{Current frontier and scope.}
The catalogue asks for a polynomial \emph{pivot count}; it does not
require polynomial-time computation of the next pivot.  That omission
is mathematically important, not a defect to be silently repaired.
\sref{30006742}{1}  The August 2026 revision of Disser--Loho--Maat--Mosis
proves lower bounds for specified ranking-based classes of pivot rules,
not for every rule allowed by this broad selection model.
\eref{30006742}{2}  Black's September 2026 preprint constructs lattice
polytopes in $[0,3]^{6d}$ with monotone diameter at least $2^d-1$.
Its parameter is geometric dimension with a coordinate bound, not
immediately the catalogue's pair $(m,n)$ for nondegenerate equality
form. \eref{30006742}{1}  These recent statements are recorded with their
hypotheses and are not treated as independently audited universal
resolutions.

\paragraph{Attack map.}
One route constructs a potential decreasing at every legal pivot and
bounded by a polynomial in $m,n$; its missing property is a uniform
range independent of numerical coefficients.  A second searches all
paths and chooses a shortest improving one, exploiting the absence
of a per-pivot efficiency condition.  That route reduces the problem
to a geometric distance question, but does not automatically bound the
distance.  A counterexample route must force \emph{every} legal path
to be long in the correct input parameters; an exponentially long
trajectory of one favorite rule is insufficient.  The shortest-path
route and its exact quantifiers are investigated here.

\paragraph{Attempt I: the directed feasible-basis graph.}
Assume redundant equations have been removed so that $A$ has row rank
$m$; this operation cannot increase either relevant input parameter.
For a bounded nondegenerate feasible rational program in standard form,
let $\mathcal B$ be its feasible bases.  Every basis corresponds to a
vertex, and nondegeneracy makes that correspondence one-to-one.  Join
two bases when a feasible single-variable exchange connects them.
Orient an edge toward strictly larger objective value; optimal bases
are terminal targets.

A nonoptimal vertex has an improving incident edge.  Here is the usual
geometric reason rather than an assumption of termination: if $v$ is
nonoptimal and $w$ is better, the feasible direction $w-v$ lies in the
polyhedral tangent cone at $v$.  That cone is generated by the incident
edge directions for a bounded polytope.  A positive value of the
objective on $w-v$ forces a positive value on at least one generator.
Nondegeneracy realizes that edge as a simplex pivot.  Every improving
walk is finite because the objective strictly increases on a finite
vertex set.  Hence each feasible basis has a directed path to an optimum.

Let $\ell(B)$ be the length of a shortest such path.  If the convention
also permits zero-objective pivots, include those edges in the legal
graph and use shortest distances there instead.  Improving paths still
ensure reachability.  The equivalence proved next works with either
specified convention; it does not exploit worsening pivots.

\textbf{Pivot-count equivalence.}
In the catalogue's model, a deterministic or randomized polynomial
pivot-count rule exists if and only if there is a polynomial $p(m,n)$
such that every allowed input and every feasible starting basis satisfy
\begin{equation}\label{30006742:distance}
 \ell(B)\leq p(m,n).
\end{equation}
Moreover, if any randomized rule has this expected bound, a uniform
deterministic rule with no worse pivot-count bound exists, although
its per-pivot running time may be exponential.

\emph{Proof.} Every realized sequence of legal pivots reaching an
optimum has length at least $\ell(B)$.  Thus an expected-length bound
implies \eqref{30006742:distance}; finite expected termination also excludes
nontermination with positive probability.  Conversely, enumerate all
$\binom nm$ candidate bases using exact rational arithmetic, test
invertibility and feasibility, and construct the legal exchange graph.
Compare the finitely many exact objective values to identify all
optimal bases.  Reverse breadth-first search from them computes
$\ell$ on every feasible basis.  Choose, with a deterministic
lexicographic tie-break, an outgoing edge from $B$ to a basis with
distance $\ell(B)-1$.  Repeating reaches an optimum in exactly
$\ell(B)$ pivots.  The enumeration is finite and uniformly computable
from the input.  It need not be efficient, and no length-dependent
advice or existence-only oracle is being used. \hfill$\square$

The construction is therefore a legitimate rule in the specified
model.  The unproved part is \eqref{30006742:distance}, not the existence
or computability of the shortest-path selector.  Enumerating all bases
without proving a short path would yield only the elementary bound
\[
 \ell(B)\leq|\mathcal B|-1\leq\binom nm-1.
\]
This is polynomial when $m$ or $n-m$ is fixed, but not for jointly
varying $m,n$.

\paragraph{Attempt II: geometric consequences and a calibration family.}
A full-dimensional simple rational $d$-polytope with $f$ facets admits
a slack-coordinate representation in $\R^f$:
\[
 s_i=b_i-a_i\cdot x\geq0\quad(1\leq i\leq f),\qquad As=b,
\]
where the affine relations among the slacks have rank $f-d$.
The affine slack map is injective for a bounded full-dimensional
polytope: a direction annihilated by all facet normals would otherwise
be a lineality direction.  A simple vertex has exactly $d$ zero slacks
and $f-d$ positive ones, giving a nondegenerate feasible basis in this
representation.  Edges and objective improvements are preserved.
Thus a polynomial bound in \eqref{30006742:distance} would imply polynomial
monotone distances to an optimum in terms of $f,d$ for such polytopes.

With a generic rational objective there is a unique optimum.  Routing
two arbitrary vertices to that same vertex would also bound their
undirected graph distance by $2p(f-d,f)$.  Hence the proposed pivot
bound has a polynomial polytope-diameter consequence.  The converse
inference is not valid: a short undirected path may decrease the
objective, and the route required by simplex is directed.  No
polynomial-diameter conjecture is assumed as an established theorem.

For an exact small-model check, consider
\[
 x_i+y_i=1,\quad x_i,y_i\geq0\quad(1\leq i\leq d),\qquad
 \max\sum_{i=1}^d w_i x_i,\quad w_i>0.
\]
Here $m=d,n=2d$, all $2^d$ vertices are nondegenerate, and an improving
pivot changes one zero $x_i$ to one.  The distance to the unique
optimum is exactly the number of zero $x_i$, at most $d$.  The number
of bases is exponential, but the shortest-path distance is linear.
This separates the cost of the brute-force selector from its pivot
count and prevents an exponential state count from being mistaken
for a geometric lower bound.

\paragraph{Attempt III: audit a plausible counterexample transfer.}
Black's recent preprint concerns genuine \emph{shortest} monotone
distances, not just one inefficient pivot trajectory: its construction
traps improving paths in an exponentially long shadow path.  This is
a relevant obstruction to any argument bounding monotone diameter by
a polynomial in geometric dimension and a fixed lattice coordinate
range. \eref{30006742}{1}, Theorem 1.1 and Section 2.

To turn it into a counterexample to the selected statement, one would
still need to give an allowed nondegenerate equality-form family with
$m,n$ polynomial in the parameter $d$, preserve the long directed
distance, and verify that the bound beats \emph{every} polynomial in
those $m,n$.  The displayed construction starts with exponentially
many listed vertices and uses a Minkowski sum.  Its geometric dimension
alone supplies no verified polynomial bound on the number of facets
in the resulting description, and a change to slack variables need
not preserve the small lattice-coordinate range.  Nor may
nondegeneracy be imposed by an arbitrary perturbation without checking
that new short improving paths have not appeared.

The preprint itself distinguishes equality-form lattice polytopes in
its introduction.  This notebook does not prove either that its family
has polynomially many facets or that such a conversion is impossible.
It records the missing conversion rather than claiming a counterexample
to a stronger or differently parametrized statement.  Likewise, the
ranking restrictions in \eref{30006742}{2} exclude the arbitrary global
shortest-path selector above, so those lower bounds do not refute
\eqref{30006742:distance}.

\paragraph{Stress test.}
Multiple optimal bases are handled by reverse search from all of them;
a tie does not require an illegal objective decrease.  The proof uses
exact rational feasibility and objective comparisons, not a numerical
tolerance that could change the directed graph.  Randomness cannot
reduce the expected number of traversed edges below shortest-path
distance.  On the other hand, the deterministic selector's arithmetic
work is unbounded by a polynomial in $m,n$, so it would not yield a
strongly polynomial-time linear-programming algorithm merely from a
polynomial pivot bound.

The companion script enumerates the complementary-variable cube bases
for small $d$, computes the improving graph and reverse distances, and
checks the distance formula against every starting basis.  It also
checks a finite directed-path example where an extra undirected edge
is unusable because it decreases the objective.  Such graph examples
are diagnostics only, not claimed realizations as arbitrary LP graphs.

\paragraph{Outcome.}
\textbf{Outcome: Conditional reduction.}

For the exact pivot-count model, existence of any expected-polynomial
rule is equivalent to a uniform polynomial shortest legal path bound,
and that bound would already admit a deterministic uniform selector.
The reduction isolates the geometric statement that must be proved or
disproved and audits the parameter and nondegeneracy obstacles in a
recent potential counterexample transfer.  No polynomial bound on
those distances, or superpolynomial lower bound in the required
parameters, has been established.  No priority claim is made.
% END RESEARCH ADDITION 131
\subsection{Sources}
\begin{itemize}\raggedright
\item[\textnormal{[S1]}]\hypertarget{source:30006742:1}{} Math Conjectures, OPT-002, catalog snapshot 31 July 2026.
\url{https://mathconjectures.com/conjectures/OPT-002}.
{\small\textit{Catalogue formulation source.}}
% BEGIN RESEARCH REFERENCES 131
\item[\textnormal{[E1]}]\hypertarget{extra:30006742:1}{} Alexander E. Black.
\emph{Monotone Diameters of Lattice Polytopes}, arXiv:2609.08647v1, 8 September 2026. Theorem 1.1 and Section 2; the geometric-dimension parameter and equality-form distinction in the introduction. Recent preprint; no conversion to the catalogue\textquotesingle s required input family is presumed.
\url{https://arxiv.org/html/2609.08647v1}.
{\small\textit{Consulted 27 September 2026 for the indicated result or scope.}}
\item[\textnormal{[E2]}]\hypertarget{extra:30006742:2}{} Yann Disser, Georg Loho, Matthew Maat, and Nils Mosis.
\emph{Lower bounds for ranking-based pivot rules}, arXiv:2512.16684v2, 10 August 2026. Abstract and introduction; restrictions on the information used by the pivot rule.
\url{https://arxiv.org/abs/2512.16684v2}.
{\small\textit{Consulted 27 September 2026 for the indicated result or scope.}}
% END RESEARCH REFERENCES 131
\end{itemize}


\end{document}
